\section*{Décider de la fenêtre et du minimum}
À minimum fixe de 1, élargir la fenêtre de $2W$ à $8W$ relève l'OTD de
35,2 points alors que la traversée moyenne reste proche de 107 h. La marge
avant échéance aide donc à absorber les aléas. À $45W$, l'en-cours a presque
doublé par rapport à $8W$ et l'OTD reste à 100\,\% sur ces graines. Une hausse
de l'en-cours seule ne suffit pas ici à dégrader le service.

Le couplage progressif donne 98,5\,\% à $8W$ et 17,2\,\% à $45W$.
L'écart en faveur du minimum fixe de 1, à fenêtre identique, passe de
1,5 point [1,0 ; 2,0] à 82,8 points [80,3 ; 85,4]. Les séries effectivement
formées grossissent, la traversée moyenne dépasse 1\,660 h et le goulot attend
des regroupements pendant 270 h à $45W$. Cette dernière mesure ne représente
qu'une partie de l'attente dans une ligne à dix postes ; elle ne permet pas
d'attribuer à elle seule toute la perte d'OTD.

Pour cette usine et cet objectif de service, un minimum de 1 et une fenêtre
de $8W$ atteignent déjà le plateau observé avec environ onze ordres en cours.
Si le lot doit croître avec la fenêtre, les données favorisent le couplage
progressif autour de $8W$ à $12W$ ; dix graines ne départagent pas ces deux
fenêtres. Le couplage rapide reste sous 14\,\% même à son meilleur point.
La branche descendante de l'OTD dépend donc de la politique adoptée. Cette
expérience quantifie un mécanisme compatible avec le raisonnement de Goldratt,
sans mesurer directement les interventions humaines.
